% Part: lambda-calculus
% Chapter: introduction
% Section: computable-lambda

\documentclass[../../../include/open-logic-section]{subfiles}

\begin{document}

\olfileid{lam}{int}{lrp}
\olsection{可计算函数是 \usetoken{S}{lambda definable}}

\begin{thm}
\ollabel{thm:computable-lambda}
每个部分可计算函数都是 $\lambd$-可定义的。
\end{thm}

\begin{proof}
我们需要证明，每个部分可计算函数~$f$ 都由某个 $\lambda$-项~$F$ 所 $\lambd$-定义。
根据 克林的范式定理，只需证明每个原始递归函数都由某个 $\lambda$-项 $\lambd$-定义，
然后证明 $\lambd$-可定义的函数对适当的复合与无界搜索封闭。
为了证明每个原始递归函数都由某个 $\lambda$-项 $\lambd$-定义，
只需证明初始函数是 $\lambd$-可定义的，
并且 $\lambd$-可定义的偏函数对复合、原始递归和无界搜索封闭。
\end{proof}

我们将使用更传统的记法，以使证明的余下部分更具可读性。
例如，我们将写 $M(x, y, z)$ 而非 $Mxyz$。
虽然这种写法具有提示性，但你应当记住，无类型 $\lambda$-演算中的项并不具有关联的元数；
因此，对于同一个项~$M$，写 $M(x,y)$ 与 $M(x,y,z,w)$ 同样是有意义的。
但使用这种记法表明我们将 $M$ 视作三元函数，有助于使定义背后的意图更加清晰。
类似地，我们会说“由 $M(x,y,z) = \dots$ 定义 $M$”，而非“由 $M =
\lambd[x][\lambd[y][\lambd[z][\dots]]]$ 定义 $M$”。

\end{document}
